% Part: turing-machines
% Chapter: machines-computations
% Section: well-behaved-machines

\documentclass[../../../include/open-logic-section]{subfiles}

\begin{document}

\olfileid{tur}{mac}{dis}
\olsection{Mesin Kang Tertib}

\begin{explain}
Ing bagean \olref[hal]{sec}, kita wis nliti mesin Turing kang nduwé
siji kahanan mandheg mligi~$h$---mesin kaya mangkono mesthi
mandheg ing kahanan~$h$, yen pancen mandheg. Kanthi cara iki,
mesin kang nduwé siji kahanan mandheg luwih ``tertib''
tinimbang mesin Turing umume kang diidinake. Ana watesan
liyane kang bisa kita tetepake marang tumindake mesin Turing.
Contone, kita uga durung nglarang mesin Turing mbusak
panandha pucuk pita ing kothak~$0$, utawa nyoba mindhah
sirah maca-nulis menyang kiwa saka kothak~$0$. (Miturut
definisi kita, ing kahanan iki sirah tetep ana ing kothak~$0$;
miturut definisi liyane, mesin mandheg.) Kadhangkala uga
perlu bisa nganggep yen mesin Turing, yen pancen mandheg,
mandheg nalika sirah maca-nulis ana ing kothak~$1$.
\end{explain}

\begin{defn}\ollabel{defn:disciplined}
Mesin Turing~$M$ diarani \emph{tertib} yen lan mung yen
\begin{enumerate}
    \item nduwé siji kahanan mandheg mligi~$h$,
    \item yen mandheg, sirah maca-nulis lagi maca kothak~$1$,
    \item ora tau mbusak simbol $\TMendtape$ ing kothak~$0$, lan
    \item ora tau nyoba mindhah sirah maca-nulis menyang kiwa saka kothak~$0$.
\end{enumerate}
\end{defn}

\begin{explain}
Kita wis ngrembug yen saben mesin Turing bisa diowahi dadi
mesin kang tumindake padha nanging nduwé kahanan mandheg
mligi. Carane mung nambah kahanan anyar~$h$, banjur nambah
instruksi $\delta(q, \sigma) = \tuple{h, \sigma, N}$ kanggo
saben pasangan $\tuple{q,\sigma}$ kang $\delta$ asale
ora ditegesi. Instruksi anyar iki mung kanggo kahanan asli; kahanan mandheg anyar ora diwenehi instruksi metu. Pancen saben mesin Turing~$M$ bisa diowahi
dadi mesin Turing tertib~$M'$ kang mandheg kanggo input
kang padha lan njaga output kang ditegesi, senajan
buktine rada dawa lan njlimet. Contone, yen mesin Turing
mandheg nalika sirah ora ana ing kothak~$1$, kita bisa
nambah sawetara instruksi supaya sirah obah menyang kiwa
nganti nemokake panandha pucuk pita, banjur obah siji
kothak menyang tengen, banjur mandheg. Kita masrahake
marang kowe kanggo mikirake carane netepi syarat liyane. Cathetan cakupan SOL6-F052: cara mulih menyang kiwa iki nganggep panandha pucuk asli isih ana lan dibedakake saka panandha ing njero input. Kanggo mesin umum kang bisa nimpa panandha, luwih dhisik gunakna simulasi kang njaga lan ngenali wates kiwa asli; mlaku kiwa wae durung cukup.
\end{explain}

\begin{ex}
    Ing \olref{fig:adder-disc}, mesin panambah saka
    \olref[una]{ex:adder} kita owahi dadi mesin tertib.
    \begin{figure}\[
\begin{tikzpicture}[->,>=stealth',shorten >=1pt,auto,node distance=2.8cm,
                    semithick]
  \tikzstyle{every state}=[fill=none,draw=black,text=black]

  \node[initial,state]         (A)              {$q_0$};
  \node[state]         (B) [right of=A] {$q_1$};
  \node[state]         (C) [below right of=B] {$q_2$};
  \node[state]         (D) [below left of=C] {$q_3$};
  \node[state]         (H) [left of=D] {$h$};

  \path (A) edge node {\TMtrans{\TMblank}{\TMstroke}{\TMstay}} (B)
           edge [loop below] node {\TMtrans{\TMstroke}{\TMstroke}{\TMright}} (A)
        (B) edge [loop below] node {\TMtrans{\TMstroke}{\TMstroke}{\TMright}} (B)
            edge node {\TMtrans{\TMblank}{\TMblank}{\TMleft}} (C)
        (C) edge node {\TMtrans{\TMstroke}{\TMblank}{\TMleft}} (D)
        (D) edge [loop above] node {\TMtrans{\TMstroke}{\TMstroke}{\TMleft}} (D)
        edge node {\TMtrans{\TMendtape}{\TMendtape}{\TMright}} (H);
\end{tikzpicture}
\]\caption{Mesin panambah kang tertib}
\ollabel{fig:adder-disc}
\end{figure}
\end{ex}

\begin{prop}\ollabel{prop:disciplined} Kanggo saben mesin Turing~$M$,
ana mesin Turing tertib~$M'$ kang mandheg kanthi output~$O$
yen $M$~mandheg kanthi output~$O$, lan ora mandheg yen
$M$~ora mandheg. Mligine, saben fungsi $f\colon\Nat^n \to \Nat$
kang bisa diitung dening mesin Turing uga bisa diitung
dening mesin Turing tertib.
\end{prop}

\begin{prob}\label{tur:mac:dis:prob:disc-succ}
Wenehana mesin tertib kang ngitung $f(x) = x+1$.
\end{prob}


\begin{prob}\label{tur:mac:dis:prob:copier}
Golekana mesin tertib kang, yen diwiwiti nganggo input $\TMstroke^n$,
ngasilake output $\TMstroke^n \concat \TMblank \concat \TMstroke^n$.
\end{prob}

\end{document}
